% 2010 年陕西卷理科第 19 题：男、女生身高频数分布直方图（样本各 40 人、30 人）。
% 来源 PDF：普通高考/2010/2010陕西理.pdf，SHA-256 8406e303e7aa1ac7813893781c903a7d9c904530028e6fd1c06062d1c32bd645。
% 原页：第 2 页（A3 三栏）第 1 栏下部，书页码 624。
% 重绘依据：照原页——左「男生」、右「女生」两幅频数直方图，横轴 身高/cm 自 160 起每 5 cm 一组
%   至 190，纵轴为频数（原点处画有断裂符号）。
%   男生各组频数依次为 2, 5, 14, 13, 4, 2（合计 40）；女生依次为 1, 7, 12, 6, 3, 1（合计 30）。
%   原页在各组上沿引出水平虚线并标注频数：男生 14、13、5、4、2；女生 12、7、6、3、1。
\begin{tikzpicture}[x=0.62cm, y=0.30cm, semithick, line join=round,
  lbl/.style={font=\scriptsize, inner sep=1pt}]

  % ================= 男生 =================
  \begin{scope}
    % 坐标轴
    \draw[->, >=stealth] (-0.55,-1.20) -- (-0.55,16.00);
    \draw[->, >=stealth] (-1.20,0) -- (7.05,0);
    \node[lbl, right] at (-0.60,14.75) {频数};
    \node[lbl, above] at (6.40,1.00) {身高/cm};

    % 原点断裂符号
    \draw[thin] (-0.98,0) -- (-0.90,1.10) -- (-0.82,-1.10) -- (-0.74,1.10) -- (-0.66,0);
    \node[lbl, below] at (-1.02,-0.20) {$0$};

    % 各组矩形
    \foreach \i/\h in {0/2, 1/5, 2/14, 3/13, 4/4, 5/2} {
      \draw (\i,0) rectangle (\i+1,\h);
    }

    % 频数虚线（自纵轴引到该频数组的左沿）
    \foreach \i/\h in {0/2, 1/5, 2/14, 3/13, 4/4} {
      \draw[thin, dashed] (-0.55,\h) -- (\i,\h);
    }
    \node[lbl, left] at (-0.62,2) {$2$};
    \node[lbl, left] at (-0.62,5) {$5$};
    \node[lbl, left] at (-0.62,4) {$4$};
    \node[lbl, left] at (-0.62,14) {$14$};
    \node[lbl, left] at (-0.62,13) {$13$};

    % 横轴刻度
    \foreach \i/\v in {0/160, 1/165, 2/170, 3/175, 4/180, 5/185, 6/190} {
      \draw[thin] (\i,0) -- (\i,-0.90);
      \node[lbl, below] at (\i,-1.05) {\v};
    }

    % 标题
    \node at (3.0,17.30) {男生};
  \end{scope}

  % ================= 女生 =================
  \begin{scope}[xshift=5.35cm]
    \draw[->, >=stealth] (-0.55,-1.20) -- (-0.55,16.00);
    \draw[->, >=stealth] (-1.20,0) -- (7.05,0);
    \node[lbl, right] at (-0.60,14.75) {频数};
    \node[lbl, above] at (6.40,1.00) {身高/cm};

    \draw[thin] (-0.98,0) -- (-0.90,1.10) -- (-0.82,-1.10) -- (-0.74,1.10) -- (-0.66,0);
    \node[lbl, below] at (-1.02,-0.20) {$0$};

    \foreach \i/\h in {0/1, 1/7, 2/12, 3/6, 4/3, 5/1} {
      \draw (\i,0) rectangle (\i+1,\h);
    }

    \foreach \i/\h in {0/1, 1/7, 2/12, 3/6, 4/3} {
      \draw[thin, dashed] (-0.55,\h) -- (\i,\h);
    }
    \node[lbl, left] at (-0.62,1) {$1$};
    \node[lbl, left] at (-0.62,7) {$7$};
    \node[lbl, left] at (-0.62,6) {$6$};
    \node[lbl, left] at (-0.62,12) {$12$};
    \node[lbl, left] at (-0.62,3) {$3$};

    \foreach \i/\v in {0/160, 1/165, 2/170, 3/175, 4/180, 5/185, 6/190} {
      \draw[thin] (\i,0) -- (\i,-0.90);
      \node[lbl, below] at (\i,-1.05) {\v};
    }

    \node at (3.0,17.30) {女生};
  \end{scope}
\end{tikzpicture}
